\subsection{TSPLIB: true gaps against published optima}\label{sec:tsplib}

Table~\ref{tab:tsplib} scores BGMA \lvl{max} and the LK-style baseline against the published optimum of
each instance. Distances are rounded as \code{EUC\_2D} prescribes. Instances up to $n=442$ were run
with ten seeds. The larger ones were run with one to three seeds, because one \lvl{max} run there costs
minutes to over an hour of CPU time. The \emph{equal-time} columns restart the LK-style search with new seeds until it has spent the median
serial cost of one BGMA \lvl{max} run on that instance. They report the best tour found within that
budget, the number of restarts it allowed, and the cumulative CPU time at which the restarts first
reached a tour at least as short as BGMA's median.


\begin{table*}[t]
\centering\scriptsize
\caption{TSPLIB \code{EUC\_2D} instances against published optima. Gaps are true gaps $L/L^\ast-1$.
``Serial time'' is the median CPU time of one \lvl{max} run (sum of its three units); with three worker
threads the wall clock is about a third of it. ``Below opt.''\ is how far the 600-iteration 1-tree bound
lies under the optimum, i.e.\ how much a bound gap would overstate the true gap.}
\label{tab:tsplib}
\setlength{\tabcolsep}{3pt}
\begin{tabular}{@{}lrrrrrrrrrrrrr@{}}
\toprule
 & & & \multicolumn{4}{c}{BGMA \lvl{max}} & \multicolumn{3}{c}{LK-style, one run} & \multicolumn{3}{c}{LK-style, equal time} & 1-tree\\
\cmidrule(lr){4-7}\cmidrule(lr){8-10}\cmidrule(lr){11-13}\cmidrule(l){14-14}
instance & $n$ & seeds & opt. & median & best & serial & opt. & median & time & best & runs & reached & below\\
 & & & hits & gap & gap & time & hits & gap & & gap & & median & opt.\\
\midrule
eil51 & 51 & 10 & 9/10 & 0.00\% & 0.00\% & 352\,ms & 1/10 & +0.23\% & 19\,ms & 0.00\% & 26 & 44\,ms & 0.82\%\\
berlin52 & 52 & 10 & 10/10 & 0.00\% & 0.00\% & 385\,ms & 10/10 & 0.00\% & 45\,ms & 0.00\% & 10 & 45\,ms & 0.00\%\\
st70 & 70 & 10 & 10/10 & 0.00\% & 0.00\% & 496\,ms & 8/10 & 0.00\% & 77\,ms & 0.00\% & 7 & 136\,ms & 0.59\%\\
kroA100 & 100 & 10 & 10/10 & 0.00\% & 0.00\% & 876\,ms & 10/10 & 0.00\% & 127\,ms & 0.00\% & 7 & 131\,ms & 1.62\%\\
pr124 & 124 & 10 & 10/10 & 0.00\% & 0.00\% & 1.44\,s & 8/10 & 0.00\% & 331\,ms & 0.00\% & 5 & 324\,ms & 1.63\%\\
bier127 & 127 & 10 & 10/10 & 0.00\% & 0.00\% & 1.58\,s & 3/10 & +0.83\% & 224\,ms & 0.00\% & 6 & 636\,ms & 0.72\%\\
ch130 & 130 & 10 & 10/10 & 0.00\% & 0.00\% & 1.68\,s & 3/10 & +0.29\% & 429\,ms & 0.00\% & 5 & 1.10\,s & 0.56\%\\
tsp225 & 225 & 10 & 10/10 & 0.00\% & 0.00\% & 7.75\,s & 3/10 & +0.37\% & 451\,ms & 0.00\% & 15 & 1.44\,s & 0.97\%\\
fl417 & 417 & 10 & 10/10 & 0.00\% & 0.00\% & 57.6\,s & 1/10 & +0.78\% & 3.05\,s & 0.00\% & 20 & 20.3\,s & 5.34\%\\
pcb442 & 442 & 10 & 0/10 & +0.66\% & +0.51\% & 64.8\,s & 0/10 & +0.56\% & 1.13\,s & +0.25\% & 60 & 1.18\,s & 0.56\%\\
u574 & 574 & 3 & 0/3 & +0.50\% & +0.33\% & 3.5\,min & 0/3 & +1.14\% & 3.23\,s & +0.47\% & 62 & 3.2\,min & 0.53\%\\
rat783 & 783 & 3 & 0/3 & +1.15\% & +1.04\% & 9.7\,min & 0/3 & +1.23\% & 1.74\,s & +0.69\% & 298 & 5.61\,s & 0.38\%\\
pr1002 & 1002 & 1 & 0/1 & +2.10\% & +2.10\% & 39.8\,min & 0/1 & +2.19\% & 5.73\,s & +1.07\% & 400 & 12.1\,s & 0.90\%\\
u1060 & 1060 & 1 & 0/1 & +1.45\% & +1.45\% & 58.4\,min & 0/1 & +2.06\% & 6.66\,s & +0.55\% & 555 & 12.8\,s & 0.86\%\\
\bottomrule
\end{tabular}

\end{table*}

\paragraph{Up to $n=442$.}
BGMA \lvl{max} returned the optimum in \val{tsSmallHits} of \val{tsSmallRuns} runs. It did so on all ten
seeds for \val{tsSmallAllSeeds} of the \val{tsSmallCount} instances and on nine seeds for eil51. It never
did on pcb442 (median \ts{pcb442}{maxMedGap}, best \ts{pcb442}{maxBestGap}). Single LK-style runs are
much less reliable: \val{tsSmallLkHits} of \val{tsSmallRuns} were optimal. They are also much cheaper,
and the best of ten LK-style runs reaches the optimum on \val{tsSmallLkAny} of the \val{tsSmallCount} instances, the same ones as BGMA. On fl417 the
ten LK-style runs together cost \ts{fl417}{lkTotal}, against \ts{fl417}{maxSerial} for one \lvl{max} run.
On pcb442 the LK-style baseline is better outright: a median of \ts{pcb442}{lkMedGap} in \ts{pcb442}{lkMs}
per run, against \ts{pcb442}{maxMedGap} for \lvl{max} in \ts{pcb442}{maxSerial}. Instance structure
matters as much as size. fl417, with its dense clusters, is solved on every seed, while pcb442, which
is only slightly larger, is solved on none.

\paragraph{Beyond $n=500$.}
The median gaps of \lvl{max} range from \val{tsBigRange}. These are true gaps: the excess is genuine tour excess,
not slack in a bound. The three portfolio units show the top of the dial flattening. On u574 the
\lvl{xhigh} unit reached \ts{u574}{armXhigh} and the two \lvl{max} restarts \ts{u574}{armMax}. On rat783
they reached \ts{rat783}{armXhigh} against \ts{rat783}{armMax}. On pr1002 they reached
\ts{pr1002}{armXhigh} against \ts{pr1002}{armMax}, and on u1060 \ts{u1060}{armXhigh} against
\ts{u1060}{armMax}. The heavier \lvl{max} configuration is not reliably better than \lvl{xhigh} on the
same instance. Its contribution is mostly that of extra independent samples. A single \lvl{max} run on
pr1002 costs \ts{pr1002}{maxSerial} of CPU time (\ts{pr1002}{maxWall} on three threads).

% filled when the equal-time data exist


\paragraph{The bound gap misleads.}
The last column of Table~\ref{tab:tsplib} shows the 1-tree bound \val{boundBelowRange} below the optimum,
depending on the instance. An optimal fl417 tour has a bound gap of \ts{fl417}{boundGap}, while an
optimal berlin52 tour has \ts{berlin52}{boundGap}. Had we reported bound gaps, as is common when no
optimum is known, fl417 would have looked like BGMA's worst instance. It is one of its best.

\subsection{Robustness, determinism and degenerate inputs}\label{sec:robust}
The stress harness ran \lvl{low} and \lvl{max} on seven families at $n=20$, $100$ and $400$
(\val{stressRuns} runs, Table~\ref{tab:stress}). Every tour was checked to be a valid permutation with
a correctly recomputed length. No run produced an invalid tour (\val{stressInvalid} failures), and no
tour came out shorter than the proven bound (\val{stressViolations} violations).

\begin{table}[t]
\centering\footnotesize
\caption{Degenerate and known-optimum families. ``Gap to bound'' is listed per $n$ ($20$, $100$, $400$)
and is a ceiling on the true gap. \emph{fine} and \emph{coarse} are the same point set scaled by
$10^{-3}$ and $10^{6}$.}
\label{tab:stress}
\setlength{\tabcolsep}{4pt}
\begin{tabular}{@{}llrrrl@{}}
\toprule
family & level & $n$ & runs & invalid & result\\
\midrule
convex & \lvl{low} & 20, 100, 400 & 3 & 0 & hull cycle (optimal) at every $n$\\
convex & \lvl{max} & 20, 100, 400 & 3 & 0 & hull cycle (optimal) at every $n$\\
line & \lvl{low} & 20, 100, 400 & 3 & 0 & known optimum at every $n$\\
line & \lvl{max} & 20, 100, 400 & 3 & 0 & known optimum at every $n$\\
grid & \lvl{low} & 20, 100, 400 & 3 & 0 & gap to bound 0.00\%, 0.00\%, +0.41\%\\
grid & \lvl{max} & 20, 100, 400 & 3 & 0 & gap to bound 0.00\%, 0.00\%, 0.00\%\\
dup & \lvl{low} & 20, 100, 400 & 3 & 0 & gap to bound 0.00\%, +1.21\%, +2.32\%\\
dup & \lvl{max} & 20, 100, 400 & 3 & 0 & gap to bound 0.00\%, +0.50\%, +0.85\%\\
same & \lvl{low} & 20, 100, 400 & 3 & 0 & known optimum at every $n$\\
same & \lvl{max} & 20, 100, 400 & 3 & 0 & known optimum at every $n$\\
fine & \lvl{low} & 20, 100, 400 & 3 & 0 & gap to bound 0.00\%, +0.95\%, +3.84\%\\
fine & \lvl{max} & 20, 100, 400 & 3 & 0 & gap to bound 0.00\%, +0.94\%, +1.65\%\\
coarse & \lvl{low} & 20, 100, 400 & 3 & 0 & gap to bound 0.00\%, +0.95\%, +3.96\%\\
coarse & \lvl{max} & 20, 100, 400 & 3 & 0 & gap to bound 0.00\%, +0.94\%, +1.65\%\\
\bottomrule
\end{tabular}

\end{table}

\paragraph{Known optima by construction.}
For points in convex position the hull cycle is optimal, and both levels returned it at every size.
At $n=1000$, \lvl{low} returned \val{convex1000@low@hull} in \val{convex1000@low} and \lvl{max}
returned \val{convex1000@max@hull} in \val{convex1000@max} of CPU time. Collinear points (optimum: twice
the span) and coincident points (optimum: 0) were solved exactly. On the $20\times 20$ grid,
\lvl{max} met the 1-tree bound, which proves its tour optimal, while \lvl{low} did not.

\paragraph{The bound can be very loose.}
On collinear points the tour is optimal by construction, yet its bound gap reads
\val{stress@line@max@100} at $n=100$ and \val{stress@line@max@400} at $n=400$. A degree-constrained relaxation cannot express the structure of a line. This
is the same phenomenon as fl417 in Section~\ref{sec:tsplib}, in a more extreme form.

\paragraph{Scale dependence.}
The \emph{fine} and \emph{coarse} families are one point set at two scales, so a scale-free method
would return the same relative gap on both. \lvl{max} nearly does (\val{stress@fine@max@400} against \val{stress@coarse@max@400} at
$n=400$). \lvl{low} does not (\val{stress@fine@low@400} against \val{stress@coarse@low@400}), because the anneal temperature is set in absolute units (Section~\ref{sec:limits}).

\paragraph{Determinism.}
Every level was run three times on each of two $n=100$ instances (\val{repeatRuns} runs in
\val{repeatGroups} groups). In \val{repeatIdentical} of the \val{repeatGroups} groups the repeats
returned the same length and the same tour. Runs are a deterministic function of (instance, seed,
level), also under the worker pool.

